\begin{textsample}{POS Dim 2 – profiled\_gpt – Score 8.00 – profiled\_gpt/t000924\_gpt.txt}  \label{ex:f2_pos_004}
well it is a bit like Chinese whispers, isn't it, but with pictures and people instead of just a sentence going round a circle.
I ’m imagining it very simply, Monday to Friday, you start with one person on the Monday, and then on the Tuesday you talk to someone they ’re connected to.
Then Wednesday ’s person is connected to Tuesday ’s, and so on through the week.
Each step along, you ’re only ever one person away from the last interview, so the chain \textbf{builds} quite naturally.
Those connections could come through social media – Facebook friends, Twitter followers, LinkedIn contacts – or just people they know in \textbf{real} life and actually meet.
As you go, you ’d start to see that they ’re not just randomly connected : they ’ve got shared problems, similar interests, similar concerns that keep popping up in slightly \textbf{different} forms.
By the time you ’ve done all five, the \textbf{set} of videos \textbf{themselves} \textbf{becomes} a kind of narrative.
You can watch them back and see this thread that runs from the original Monday person, through each link, right to the Friday person.
Each interview doesn't just sit on its \textbf{own} ; it pushes the story forward, adds a bit \textbf{more} to that picture of what connects them all, and in the end everything sort of loops back conceptually to where you began.

% matched lemmas: become, build, different, more, own, real, set, themselves
\end{textsample}
